% Insert this subsection immediately after Table~\ref{tab:llm-judgment-results}
% and before \section{Results analysis and limitations}.

\subsection{Exploratory source-level analysis}

Because the reference sample contains ten documents from each source, the same
annotations can be inspected by source to identify possible sensitivity to
documentary format. This analysis is exploratory: ten documents per source are
not sufficient to estimate a stable source-specific performance level, and the
results should therefore be read as diagnostic evidence rather than as a ranking
of sources.

Target audience, initiative type, and theme were fully accepted in all six
source subsets. Territorial scale was also fully correct for five sources. The
only exception was Fondation de France, for which one expected scale label was
missed, resulting in an F1 score of 0.957. Table~\ref{tab:llm-source-results}
focuses on the fields for which source-dependent variation is most informative.
For territory, holder, and date, the notation $a/b$ gives the number of accepted
predictions over evaluable documents; ``--'' indicates that no document was
evaluable for that field.

\begin{table}[H]
\centering
\small
\begin{tabular}{lrrrr}
\toprule
Source & Scale F1 & Territory & Holder & Date \\
\midrule
CNAV & 1.000 & 8/9 (0.889) & 10/10 (1.000) & 4/8 (0.500) \\
Fiches actions SIAGE & 1.000 & 10/10 (1.000) & 8/9 (0.889) & 8/10 (0.800) \\
Fondation de France & 0.957 & 8/8 (1.000) & 7/7 (1.000) & 8/8 (1.000) \\
HCFEA & 1.000 & 9/9 (1.000) & 2/2 (1.000) & 9/10 (0.900) \\
IReSP Autonomie & 1.000 & 8/8 (1.000) & -- & 9/10 (0.900) \\
RFVAA & 1.000 & 10/10 (1.000) & 10/10 (1.000) & 6/10 (0.600) \\
\bottomrule
\end{tabular}
\caption{Exploratory source-level results from the stratified reference sample.
Territory, holder, and date are reported as accepted predictions over evaluable
documents, followed by the accepted rate.}
\label{tab:llm-source-results}
\end{table}

The results suggest that date extraction is particularly dependent on source
format. The accepted date rate ranges from 0.500 for CNAV to 1.000 for
Fondation de France, with RFVAA also showing a lower rate of 0.600. This is
consistent with the different roles of temporal expressions in the sources:
they may describe an initiative, an event, a publication, or an update to a
document. The IReSP holder field was non-evaluable in all ten sampled records,
which reflects the frequent absence of a clearly identifiable initiative holder
in research-project descriptions rather than a measured model failure.
